\documentclass[11pt]{letter}
\usepackage[margin=1.1in]{geometry}
\signature{Krishna Harish\\Elkins High School, Missouri City, TX\\krishnaharish2009@gmail.com}
\address{Krishna Harish\\Elkins High School\\Missouri City, TX, USA}
\begin{document}
\begin{letter}{The Editors\\Journal of Applied and Computational Topology}

\opening{Dear Editors,}

Please consider the enclosed manuscript, ``Directed Weisfeiler--Leman
refinement versus the invariants of the magnitude--path spectral
sequence,'' for publication in the Journal of Applied and Computational
Topology.

The magnitude--path spectral sequence organizes magnitude homology, path
homology, and reachability homology of a digraph into pages of a single
object (Asao 2023; Hepworth and Roff 2024, 2025), while directed
Weisfeiler--Leman color refinement bounds the separating power of
message-passing neural networks on digraphs. For undirected graphs the
analogous comparison, persistent homology against the WL hierarchy, was
carried out by Ballester and Rieck (LoG 2024); no placement of the
directed homology theories relative to directed refinement existed. This
paper carries out that comparison exhaustively and exactly over all
1,550,786 digraphs on three to six vertices, and proves, among other
results: the two families are incomparable, with minimal witnesses
determined exactly; the bigraded rank function of the first page of the
spectral sequence determines neither its second page nor its target
(minimal witnesses on five vertices, via a family in which the first
page degenerates to counting while path homology recovers the Betti
numbers of the underlying graph); magnitude homology separates digraphs
on which both reciprocity-aware refinement and the directed distance
multiset agree; and a component-counting description of degree-two
magnitude homology, building on a basis of Ivanov and Mukoseev,
identifies the mechanism (round-trip structure modulo geodesic
fillability) with closed formulas.

Every claim is either proved or is a finite computation in exact
rational arithmetic; two independent implementations cross-validate
each other, the degree-two formulas were verified mechanically on all
1,550,786 digraphs, and every witness pair was re-verified from
scratch. All code and data are public at
https://github.com/krishnatheaverage/mpss-dwl-expressivity (release
v1.0.0), including scripts reproducing every number and table in the
paper.

This work is not under consideration elsewhere, and I have no competing
interests. Researchers active on the magnitude--path spectral sequence
and on the expressivity of topological descriptors in graph learning,
for example the author groups cited above, would be well placed to
referee it.

\closing{Sincerely,}
\end{letter}
\end{document}
